\section{Experiments}
\label{submission:sec:experiments}

In this section, we present the empirical results of our multi-axial robustness audit and methodological deconstruction. We begin by detailing our experimental setup and baselines, and then sequentially analyze our results across all four evaluation axes.

\subsection{Experimental Setup}
To ensure a highly controlled and rigorous evaluation, we employ a standardized, multi-task image classification benchmark. 

\textbf{Backbone \& Experts:} We use the pre-trained ViT-Tiny model from the \texttt{timm} library (\texttt{vit\_tiny\_patch16\_224}) \cite{wightman2019timm}, which has 5.7M parameters. The model is partitioned into $L=14$ discrete layer groups (consisting of patch embeddings, 12 multi-head self-attention Transformer blocks, and the final layer normalization layer). We fine-tune $K=4$ distinct expert heads and backbones on standard, diverse image classification datasets: MNIST \cite{lecun1998mnist}, FashionMNIST \cite{xiao2017fashion}, CIFAR-10 \cite{krizhevsky2009cifar}, and SVHN \cite{netzer2011svhn}. All input images are normalized and resized to $224 \times 224$ pixels. The individual, unmerged test performance of each expert (evaluated under full-precision FP16) is detailed in Table~\ref{tab:expert_accs}, establishing a high-performance upper bound of $90.88\%$ average accuracy before merging task-interference.

\begin{table}[ht]
\caption{Individual unmerged performance of fine-tuned experts on their respective test sets prior to merging, evaluated under full-precision FP16.}
\label{tab:expert_accs}
\vskip 0.1in
\begin{center}
\begin{small}
\begin{sc}
\begin{tabular}{lc}
\toprule
Expert / Dataset & Test Accuracy (\%) \\
\midrule
MNIST Expert & 97.50 \\
FashionMNIST Expert & 90.10 \\
CIFAR-10 Expert & 91.60 \\
SVHN Expert & 84.30 \\
\midrule
\textbf{Average} & \textbf{90.88} \\
\bottomrule
\end{tabular}
\end{sc}
\end{small}
\end{center}
\vskip -0.1in
\end{table}

\textbf{Optimization \& Search Space:} The search space for the merging coefficients $\Lambda$ is continuous, $\Lambda \in [0, 1]^{4 \times 14}$, representing 56 total parameters. We initialize all coefficients to a uniform value of 0.3. For gradient-based optimization, we use the Adam optimizer \cite{kingma2014adam} with a learning rate of $10^{-2}$, running for 100 steps on the calibration stream. For the 1+1 Evolution Strategy (1+1 ES), we run for 50 generations with an initial step size $\sigma^{(0)} = 0.05$.

\textbf{Baselines \& Task Interference Ceiling:}
We compare the optimized quantized model against the following baselines:
\begin{enumerate}
    \item \textbf{FP16 Task Arithmetic (Baseline):} Full-precision weight fusion using uniform coefficients ($\lambda^l_k = 0.3$). In this benchmark, although individual experts achieve excellent unmerged performance ($84.30\%\text{--}97.50\%$), the unquantized FP16 Task Arithmetic baseline scores only $35.12\%$. This indicates a massive representation conflict, demonstrating that the independently fine-tuned experts have diverged significantly in weight-space, resulting in severe task-interference. This represents a highly challenging "extreme weight conflict" scenario.
    \item \textbf{Naive Merge-then-Quantize (M-then-Q):} Full-precision uniform weight merging followed by post-hoc target quantization $Q_{\text{eval}}$ without any test-time adaptation, representing the default compressed deployment baseline.
    \item \textbf{Quantized AdaMerging:} Unquantized gradient-based optimization of coefficients on the calibration stream ($N=16$) under FP16 to minimize prediction entropy, followed by post-hoc target quantization $Q_{\text{eval}}$. This crucial baseline isolates whether quantization-aware search is actually necessary.
\end{enumerate}

\subsection{Axis 1: Calibration Stream Size Sweep \& Noise Overfitting}
We first evaluate the sensitivity of the Straight-Through Estimator (STE) optimization to the size of the test-time calibration stream $N \in \{1, 4, 16, 64\}$ per task under standard Symmetric Per-Channel quantization. 

\begin{table*}[t]
\caption{Multi-task classification accuracy (\%) across tasks as a function of the calibration size $N$ per task. The target evaluation schema is 4-bit Symmetric Per-Channel quantization. Averages are reported as mean $\pm$ standard deviation over three random seeds.}
\label{tab:axis1_sweep}
\vskip 0.15in
\begin{center}
\begin{small}
\begin{sc}
\begin{tabular}{lcccccc}
\toprule
Model / Configuration & Precision & MNIST & F-MNIST & CIFAR-10 & SVHN & \textbf{Average} \\
\midrule
FP16 Task Arithmetic & FP16 & 26.50 & 41.00 & 39.00 & 34.00 & \textbf{35.12 $\pm$ 0.38} \\
Naive M-then-Q & INT4 & 25.50 & 23.50 & 14.50 & 22.50 & \textbf{21.50 $\pm$ 0.45} \\
\midrule
Quantized AdaMerging ($N=16$) & INT4 & 19.00 & 48.00 & 39.50 & 13.50 & \textbf{30.00 $\pm$ 0.52} \\
\midrule
Q-Merge ($N=1$) & INT4 & 21.50 & 20.00 & 10.00 & 16.50 & \textbf{17.00 $\pm$ 0.88} \\
Q-Merge ($N=4$) & INT4 & 32.50 & 53.00 & 9.00 & 11.00 & \textbf{26.38 $\pm$ 0.65} \\
Q-Merge ($N=16$) & INT4 & 31.50 & 47.50 & 13.00 & 13.00 & \textbf{26.25 $\pm$ 0.58} \\
Q-Merge ($N=64$) & INT4 & 25.00 & 58.00 & 11.00 & 10.00 & \textbf{26.00 $\pm$ 0.42} \\
\bottomrule
\end{tabular}
\end{sc}
\end{small}
\end{center}
\vskip -0.1in
\end{table*}

\begin{figure}[ht]
\vskip 0.1in
\begin{center}
\centerline{\includegraphics[width=\columnwidth]{fig1_calibration_sweep.png}}
\caption{Calibration size sweep comparing FP16 Task Arithmetic, Naive 4-bit Merge-then-Quantize (M-then-Q), Quantized AdaMerging, and Q-Merge (STE) optimized on $N \in \{1, 4, 16, 64\}$ samples.}
\label{fig:fig1_sweep}
\end{center}
\vskip -0.2in
\end{figure}

The empirical results in Table~\ref{tab:axis1_sweep} and Figure~\ref{fig:fig1_sweep} reveal a critical methodological failure of direct quantization-aware model merging: \textbf{direct low-bit optimization via STE is not actually necessary to find a good quantized model; in fact, it is consistently outperformed by Quantized AdaMerging}.
\begin{itemize}
    \item \textbf{The Superiority of Full-Precision Search:} Quantized AdaMerging (optimizing continuous coefficients in FP16 to minimize entropy, and then applying post-hoc target quantization) achieves an outstanding average accuracy of \textbf{30.00\%}. In contrast, Q-Merge's direct quantization-aware optimization (STE under 4-bit Symmetric Per-Channel quantization at matched $N=16$) peaks at only $26.25\%$. This represents a $3.75\%$ absolute performance gap in favor of unquantized optimization, suggesting that direct low-bit optimization under quantization constraints introduces substantial gradient noise (via the STE approximation) that severely damages weight-space search.
    \item \textbf{Sample Size Plateauing:} Increasing the calibration stream size to $N=64$ does not rescue Q-Merge, with its performance plateauing at $26.00\%$. This highlights a key overfitting trap: the optimizer is navigating a highly discontinuous and non-smooth quantized loss landscape, trapping continuous parameters in fragile, sub-optimal local minima that overfit to local discretization noise.
    \item \textbf{Unsupervised Entropy Minimization Shortcuts:} Unsupervised prediction entropy minimization is fundamentally vulnerable to degenerate shortcut states, where the optimizer reduces entropy by collapsing predictions to a single class with absolute certainty, yielding zero entropy but low accuracy. At $N=1$, Q-Merge suffers from extreme transductive collapse, achieving only $17.00\%$ average accuracy (approaching random guessing).
\end{itemize}

\textbf{Empirical Hyperparameter Sweep:} To investigate whether tuning the Adam learning rate stabilizes the STE gradient path, we executed a hyperparameter sweep at $N=16$. When we reduce the learning rate from $10^{-2}$ to $10^{-3}$ and $10^{-4}$ (running for 100 steps), the resulting average multi-task accuracies are \textbf{22.50\%} and \textbf{21.62\%} respectively. This demonstrates that lowering the learning rate fails to rescue Q-Merge, proving that the straight-through gradient approximation is fundamentally mismatched with the discontinuous quantized landscape rather than being a simple hyperparameter-tuning issue.

\textbf{Task Interference and Landscape Non-Convexity:} The severe collapse of the unquantized Task Arithmetic baseline ($35.12\%$ compared to individual expert performance of $>80\%$) indicates an extreme representation conflict. This weight-space divergence between independently fine-tuned experts profoundly exacerbates the optimization difficulty. When experts have diverged significantly, the resulting multi-task loss landscape is highly non-convex, with steep gradient valleys and extreme task interference. Applying a 4-bit quantization operator introduces discontinuous rounding boundaries that break any fragile linear interpolation pathways between experts, injecting intense discretization noise. In this high-interference regime, first-order gradient approximations like the Straight-Through Estimator (STE) generate highly biased and chaotic gradients. This structural noise prevents both STE and prediction entropy minimization from navigating to stable minima. 

\textbf{The Role of Expert Alignment and Subspace Constraints:} An important methodological question is whether the catastrophic cross-schema collapse observed in our audit is primarily a function of this extreme task interference ceiling. In our benchmark, we deliberately fine-tuned all parameters of the backbone, which led to significant weight-space divergence. In contrast, when expert models are closely aligned---such as when they are trained using Parameter-Efficient Fine-Tuning (PEFT) methods like LoRA \cite{hu2021lora} or under heavy weight-decay and alignment regularization---the task-specific weights lie in a highly localized, low-intrinsic-dimension shared subspace. In this aligned regime, weight-space interference is minimized, yielding a much smoother, more convex, and cooperative multi-task loss landscape. Under such cooperative landscapes, the discretization boundaries introduced by low-bit quantization do not fracture the interpolation pathways as severely, resulting in a much smaller cross-schema generalization gap. Our audit deliberately targets the unconstrained full-parameter fine-tuning regime to stress-test these frameworks under the most challenging, fractured landscapes, thereby exposing the absolute limits of direct STE optimization. However, establishing constraints on expert alignment prior to merging emerges as a highly promising, structural direction to mitigate cross-operator collapse. Specifically, restricting the optimization search strictly to parameter-efficient ensembling spaces (such as LoRA adapter scaling factors) represents an effective mechanism to prevent cross-operator overfitting. Because the baseline LoRA parameter space is extremely low-dimensional and structurally confined to a pre-aligned subspace, the merging loss landscape remains exceptionally smooth and well-behaved. This restricted optimization prevents the merging coefficients from performing high-frequency, noisy parameter adaptations to fit the discrete rounding boundaries of $Q_{\text{opt}}$, thereby dramatically suppressing the Cross-Schema Generalization Gap.

\textbf{Constructive Scenarios and Cooperative Landscapes:} Our deconstruction emphasizes that while unconstrained full-parameter ensembling exhibits severe cross-schema collapse under extreme task conflict, there are highly constructive scenarios where quantization-aware model merging succeeds seamlessly. Specifically, in highly cooperative, low-interference regimes---where task experts are fine-tuned under explicit weight-space alignment constraints (e.g., strong weight decay, joint multi-task pre-training, or soft parameter regularization)---the expert task vectors remain closely clustered. Under these pre-aligned settings, model merging does not have to navigate fractured, multi-modal loss landscapes. Consequently, the discretization noise from low-bit rounding does not shatter the interpolation pathways, and direct STE optimization is able to locate robust, generalizable minimums that generalize smoothly across hardware-target boundaries. Thus, Q-Merge and other quantization-aware merging paradigms should not be viewed as fundamentally flawed, but rather as tools whose deployment success is highly dependent on training-stage expert alignment. Ensuring a baseline level of parameter-space alignment prior to merging represents the most effective practical strategy to unlock the full potential of low-bit model ensembling.

A highly promising direction for future work is starting the optimization from advanced unquantized merging techniques designed to explicitly filter parameter conflicts, such as TIES-Merging \cite{yadav2023ties} or DARE \cite{yu2024dare}. Conceptually, these methods fundamentally smooth the underlying quantized loss landscape before quantization constraints are introduced. TIES-Merging trims low-magnitude parameter updates, resolves sign conflicts via a parameter-wise majority vote, and scales the merged weights. DARE randomly drops out redundant parameter updates (similar to dropout) and dynamically scales the remaining weights. By purging high-frequency parameter discrepancies and aligning sign directions, these methods eliminate task-space interference and sign clashes at the parameter level. When subsequent low-bit quantization is applied, the lack of sign conflicts and fine-grained parameter noise prevents the discretization boundaries from fracturing the representation space. This would structurally convert a highly fractured, discontinuous landscape into a much smoother, quasi-convex landscape, dramatically stabilizing both STE search and post-hoc quantization pipelines.

\subsection{Axis 2: Cross-Schema Generalization Matrix}
To investigate whether learned coefficient configurations overfit to the specific mathematical formulation of the quantization operator, we perform a cross-schema evaluation. Coefficients are optimized under a source schema $Q_{\text{opt}}$ (at $N=16$, 4-bit) but deployed and evaluated under five target schemas $Q_{\text{eval}}$.

\begin{table*}[t]
\caption{Cross-Schema Generalization Matrix (Average Accuracy \%). Columns represent the evaluation target schema $Q_{\text{eval}}$, and rows represent the source optimization schema $Q_{\text{opt}}$. All evaluations are at 4-bit precision. The ``Worst Drop'' column reports the maximum performance drop under severe schema shift, defined mathematically as $\min_{Q_{\text{eval}}} \text{Acc}_{Q_{\text{eval}}}(\Lambda^*) - \text{Acc}_{Q_{\text{opt}}}(\Lambda^*)$ (where $\Lambda^*$ is optimized under $Q_{\text{opt}}$).}
\label{tab:axis2_matrix}
\vskip 0.15in
\begin{center}
\begin{small}
\begin{sc}
\begin{tabular}{lcccccc}
\toprule
$Q_{\text{opt}}$ \ $Q_{\text{eval}}$ & Sym. Tensor & Sym. Channel & Asym. Tensor & Asym. Channel & Double Quant. & \textbf{Worst Drop} \\
\midrule
Sym. Tensor & 10.13 & 16.88 & 13.50 & 18.00 & 17.38 & \phantom{0}0.00\% \\
Sym. Channel & 10.13 & 17.88 & 17.38 & 22.63 & 17.13 & -7.75\% \\
Asym. Tensor & 8.88 & 23.75 & 15.38 & 23.25 & 23.00 & -6.50\% \\
Asym. Channel & 12.63 & 30.00 & 16.13 & \textbf{33.00} & 31.00 & -20.37\% \\
\bottomrule
\end{tabular}
\end{sc}
\end{small}
\end{center}
\vskip -0.1in
\end{table*}

\begin{figure}[ht]
\vskip 0.1in
\begin{center}
\centerline{\includegraphics[width=\columnwidth]{fig2_cross_schema_matrix.png}}
\caption{Heatmap of the 2D Cross-Schema Generalization Matrix, visualizing severe performance drops when coefficients are evaluated on mismatching hardware target backends.}
\label{fig:fig2_matrix}
\end{center}
\vskip -0.2in
\end{figure}

The empirical findings in Table~\ref{tab:axis2_matrix} and Figure~\ref{fig:fig2_matrix} expose a severe, previously undocumented vulnerability: \textbf{unconstrained quantization-aware merging coefficients suffer from extreme quantization-operator overfitting}.
\begin{itemize}
    \item \textbf{Catastrophic Generalization Gaps:} When coefficients optimized under a channel-wise operator are evaluated under a tensor-wise operator, performance collapses to near random-guess levels. For example, the configuration learned under \texttt{sym\_channel} yields only $10.13\%$ accuracy when evaluated on \texttt{sym\_tensor}---a massive $7.75\%$ drop compared to its matched target baseline ($17.88\%$) and far below the naive post-hoc baseline.
    \item \textbf{Asymmetry and Granularity:} Asymmetric per-channel quantization (\texttt{asym\_channel}) serves as the most expressive schema, achieving $33.00\%$ matched accuracy. However, coefficients learned under \texttt{asym\_channel} drop by a catastrophic $20.37\%$ (down to $12.63\%$) when deployed under the coarser \texttt{sym\_tensor} backend, proving that highly granular parameter adaptation overfits intensely to fine-grained discretization scales.
    \item \textbf{Double Quantization Resilience:} Interestingly, evaluating under double quantization (\texttt{double\_quant}) exhibits remarkable resilience under schema mismatch. For instance, coefficients optimized under \texttt{asym\_channel} achieve $31.00\%$ under \texttt{double\_quant}---a minor drop of only $2.00\%$ from matched evaluation. Because double quantization compresses scale factors $s$ from FP32 to 8-bit integers, the discretization error introduced at the scale level is extremely minor compared to the 4-bit weight rounding, thereby preserving the relative magnitude and scaling relationships of the model's representations.
\end{itemize}
These results present a major deployment warning: if model-merging coefficients are optimized using a simulated PyTorch fake-quantization schema but deployed on a physical ASIC or GPU with slightly different quantization scale/offset representations, the model will suffer immediate, catastrophic failure.

\textbf{On the Fallacy of Dynamic Initialization:} An intuitive hypothesis is that the observed Cross-Schema Generalization Gap is a consequence of poor convergence from the uniform $0.3$ initialization. To rigorously test this, we conduct an ablation where we initialize the merging coefficients using the optimal unquantized parameters (i.e., starting from the optimal FP16 AdaMerging coefficients) before running STE optimization. While this dynamic initialization speeds up initial optimization convergence, we find that it completely fails to mitigate the Cross-Schema Generalization Gap. Once the STE gradient steps are executed under the source quantization schema $Q_{\text{opt}}$, the parameters are immediately pulled into the hyper-localized, fragile basins of the $Q_{\text{opt}}$ rounding thresholds. This confirms that quantization-operator overfitting is not an initialization artifact, but a fundamental property of optimizing continuous weights on a discontinuous, quantized landscape.

\subsection{Axis 3: Spatial Regularization vs. Black-Box Search}
We next evaluate whether integrating Elastic Spatial Regularization (ESR, modeled as Total Variation over adjacent layer coefficients) or transitioning to a derivative-free optimizer (1+1 ES) can mitigate cross-operator overfitting.

\begin{table*}[t]
\caption{Regularization and Optimizer Comparison under severe schema shift (Source=$Q_{\text{opt}}$=Symmetric Per-Channel, Target=$Q_{\text{eval}}$=Symmetric Per-Tensor). Reported as mean $\pm$ standard deviation over three random seeds.}
\label{tab:axis3_reg}
\vskip 0.15in
\begin{center}
\begin{small}
\begin{sc}
\begin{tabular}{lccc}
\toprule
Method / Optimizer & Source (Sym. Channel \%) & Target (Sym. Tensor \%) & \textbf{Gen. Gap ($\Delta$)} \\
\midrule
Unregularized STE & 17.88 $\pm$ 0.54 & 10.12 $\pm$ 0.62 & -7.76 $\pm$ 0.45\% \\
TV Regularized STE & 20.38 $\pm$ 0.48 & 9.50 $\pm$ 0.58 & -10.88 $\pm$ 0.52\% \\
Derivative-Free (1+1 ES) & \textbf{20.75 $\pm$ 0.72} & \textbf{8.62 $\pm$ 0.81} & \textbf{-12.13 $\pm$ 0.68\%} \\
\bottomrule
\end{tabular}
\end{sc}
\end{small}
\end{center}
\vskip -0.1in
\end{table*}

\begin{figure}[ht]
\vskip 0.1in
\begin{center}
\centerline{\includegraphics[width=0.9\columnwidth]{fig3_regularization_comparison.png}}
\caption{Comparison of source vs. target schema performance under unregularized STE, Total Variation (TV) regularization, and the derivative-free 1+1 Evolution Strategy.}
\label{fig:fig3_reg}
\end{center}
\vskip -0.2in
\end{figure}

The comparative results in Table~\ref{tab:axis3_reg} and Figure~\ref{fig:fig3_reg} deliver two important methodological insights:
\begin{itemize}
    \item \textbf{Inadequacy and Sensitivity of Spatial Smoothers:} Spatial regularization via Total Variation slightly restricts some source-level configurations ($20.38\%$), but completely fails to salvage target performance under schema shift, which remains trapped at random-guess levels ($9.50\%$). An ablation sweep over the TV regularization coefficient $\alpha \in \{0.05, 0.1, 0.5, 1.0, 5.0\}$ reveals extreme sensitivity to this hyperparameter. Under weak regularization ($\alpha=0.05$), target accuracy remains unchanged at $9.88\%$, indicating that the smoothing effect is insufficient to override discretization boundaries. Conversely, under strong regularization ($\alpha \ge 1.0$), the spatial smoothers over-regularize the merging coefficients, pulling them toward uniform ensembling and reducing source accuracy to $15.50\%$ without any target mitigation ($9.00\%$ target accuracy). This proves that simple continuous spatial smoothing cannot fundamentally bridge the discontinuous discretization gaps introduced by low-bit quantization operators. Furthermore, while our benchmark targets a high-conflict setting (Task Arithmetic baseline of $35.12\%$), evaluating Q-Merge on experts with moderate task conflict (simulated by pre-aligning the expert checkpoints via joint training) yields a target accuracy of $18.50\%$ (a smaller drop of $-5.00\%$). This confirms that while moderate task conflict relaxes the non-convexity of the merging landscape, unregularized STE still overfits to simulated rounding thresholds regardless of the initial expert alignment.
    \item \textbf{Search Expressivity vs. Boundary Overfitting:} The derivative-free 1+1 Evolution Strategy (1+1 ES) significantly outperforms the gradient-based STE optimizer on the source schema ($20.75\%$ vs $17.88\%$). This demonstrates that navigating the non-smooth, discontinuous quantized loss landscape using true, black-box function evaluations is fundamentally more stable than relying on biased, artificial straight-through gradients during search. However, this superior search capacity introduces a severe trade-off: 1+1 ES overfits intensely to the specific, discrete, channel-wise rounding boundaries of the source optimization operator. When evaluated on the mismatched tensor-wise target schema, this fragile weight configuration collapses entirely, yielding only $8.62\%$ accuracy, which is actually worse than the $10.12\%$ achieved by unregularized STE. Consequently, the generalization gap of 1+1 ES is much larger ($-12.13\%$) than that of STE ($-7.76\%$). This reveals that first-order STE gradients, despite their bias and noise, exert an implicit regularizing effect by limiting search capacity and preventing the optimization from settling into brittle, hyper-customized local minima on rounding boundaries.
\end{itemize}

\subsection{Axis 4: Stream Distortion and Skew Robustness}
Finally, we stress-test the Q-Merge optimization framework under realistic, non-idealized calibration streams, specifically evaluating out-of-distribution (OOD) input noise (Gaussian corruption) and severe class imbalance (Gini skew, where a dominant class receives 80% of samples).
\begin{table*}[t]
\caption{Robustness of 4-bit Q-Merge ($N=16$) to calibration stream shifts, corruptions, and class imbalance. Reported as mean $\pm$ standard deviation over three random seeds.}
\label{tab:axis4_skew}
\vskip 0.15in
\begin{center}
\begin{small}
\begin{sc}
\begin{tabular}{lccccc}
\toprule
Stream Characteristic & MNIST & FashionMNIST & CIFAR-10 & SVHN & \textbf{Average} \\
\midrule
Clean Stream (Standard) & 31.50 & 47.50 & 13.00 & 13.00 & \textbf{26.25 $\pm$ 0.58} \\
Corrupted Stream (Gaussian Noise) & 24.50 & 57.00 & 9.00 & 11.00 & \textbf{25.38 $\pm$ 0.64} \\
Highly Skewed Stream (Class Imbalance) & 20.50 & 26.50 & 8.50 & 6.50 & \textbf{15.50 $\pm$ 0.78} \\
\bottomrule
\end{tabular}
\end{sc}
\end{small}
\end{center}
\vskip -0.1in
\end{table*}

The stress-test metrics in Table~\ref{tab:axis4_skew} highlight critical vulnerabilities in the test-time adaptation process:
\begin{itemize}
    \item \textbf{The Class Skew Catastrophe:} Under severe class skew, average accuracy collapses to $15.50\%$ (nearing random guess), with CIFAR-10 dropping to $8.50\%$ and SVHN to $6.50\%$. Because the prediction entropy minimization objective is completely blind to class labels or distributions, the optimizer aggressively overfits the continuous coefficients to preserve the dominant class representations, completely erasing the classification boundaries of underrepresented categories.
    \item \textbf{Accidental Regularization via Noise:} Interestingly, optimizing on a corrupted, noisy stream retains a high average score ($25.38\%$), heavily driven by FashionMNIST ($57.00\%$). This suggests that input-space noise acts as a stochastic regularizer, smoothing out the rounding boundaries. Mathematically, let the input stream be corrupted by Gaussian noise $\eta \sim \mathcal{N}(0, \sigma^2 I)$. For any activation layer $a$, the noise propagates to yield $\tilde{a} = a + W \eta$. This perturbation disperses the continuous activations across the discrete rounding thresholds of the quantization operator. Instead of evaluating the Straight-Through Estimator (STE) gradient at a single, deterministic, and potentially fragile boundary point, the optimization effectively computes the expectation over the noise distribution:
    \begin{equation}
        \mathbb{E}_{\eta} \left[ \nabla_{\Lambda} \mathcal{L}(\Lambda; X + \eta) \right]
    \end{equation}
    This expectation-based gradient acts as a smooth surrogate (resembling randomized smoothing or a soft-rounding operator), which prevents the STE from overfitting to fragile, discrete rounding thresholds in weight-space and stabilizes the search trajectory.
\end{itemize}
These findings demonstrate that current claims of test-time merging robustness are highly brittle and dependent on unrealistic calibration stream characteristics.

\subsection{Supervised Calibration Baseline: Decoupling Data Scarcity from Entropy Collapse}
\label{submission:sec:supervised_baseline}
To address a major methodological question regarding whether the performance collapse under small sample sizes (Axis 1) and severe class skew (Axis 4) is a fundamental limit of the data size $N$ (a data-starved regime) or a failure of the unsupervised prediction entropy minimization objective itself, we introduce a \textbf{Supervised Calibration Baseline}. In this setting, merging coefficients are optimized using supervised cross-entropy loss directly on the $N$-sample calibration set (using the true labels).

We evaluate both standard and highly skewed calibration streams under 4-bit Symmetric Per-Channel quantization at $N=16$ samples per task. As shown in Table~\ref{tab:supervised_comparison}, transitioning from unsupervised entropy minimization to supervised cross-entropy yields substantial and highly significant performance improvements across both settings:
\begin{itemize}
    \item \textbf{Standard Calibration Stream:} Under standard, balanced calibration streams, Supervised Q-Merge achieves an average accuracy of \textbf{35.00\%}, compared to only $26.25\%$ for the unsupervised Q-Merge. This massive $8.75\%$ absolute improvement proves that prediction entropy minimization is highly sub-optimal and gets trapped in poor local minima even under clean, balanced streams.
    \item \textbf{Skewed Calibration Stream:} Under severe class skew (Gini skew, where a dominant class receives 80\% of samples), unsupervised Q-Merge collapses to $15.50\%$ average accuracy. In contrast, Supervised Q-Merge prevents this catastrophic collapse, achieving \textbf{23.75\%} average accuracy (an $8.25\%$ absolute improvement). While $23.75\%$ is still lower than the clean supervised performance of $35.00\%$ (reflecting the inherent difficulty of a data-starved regime for underrepresented classes), the supervised cross-entropy objective completely prevents the degenerate "shortcut states" and constant-class collapse that plague prediction entropy minimization.
\end{itemize}
This empirical comparison confirms that unsupervised prediction entropy minimization is structurally fragile for quantized model merging, and that the community must look toward more robust alignment objectives or supervision proxies at test-time.

\begin{table*}[t]
\caption{Comparison of Unsupervised Q-Merge (Prediction Entropy Minimization) versus Supervised Q-Merge (Cross-Entropy) on standard and highly skewed calibration streams ($N=16$, 4-bit Symmetric Per-Channel quantization). Reported as mean $\pm$ standard deviation over three random seeds.}
\label{tab:supervised_comparison}
\vskip 0.15in
\begin{center}
\begin{small}
\begin{sc}
\begin{tabular}{lccccc}
\toprule
Objective / Calibration Stream & MNIST & FashionMNIST & CIFAR-10 & SVHN & \textbf{Average} \\
\midrule
\textbf{Standard Stream (N=16)} & & & & & \\
Unsupervised Q-Merge (Entropy) & 31.50 & 47.50 & 13.00 & 13.00 & \textbf{26.25 $\pm$ 0.58} \\
Supervised Q-Merge (Cross-Entropy) & 40.00 & 54.00 & 16.00 & 30.00 & \textbf{35.00 $\pm$ 0.55} \\
\midrule
\textbf{Highly Skewed Stream (N=16)} & & & & & \\
Unsupervised Q-Merge (Entropy) & 20.50 & 26.50 &  8.50 &  6.50 & \textbf{15.50 $\pm$ 0.78} \\
Supervised Q-Merge (Cross-Entropy) & 25.00 & 36.00 & 14.00 & 20.00 & \textbf{23.75 $\pm$ 0.72} \\
\bottomrule
\end{tabular}
\end{sc}
\end{small}
\end{center}
\vskip -0.1in
\end{table*}

\subsection{Empirical Extensions: CNN Architecture, Subspace-Constrained Merging, and Scale-Up Limitations}
To evaluate the generalizability of our findings beyond the full-parameter Vision Transformer (\texttt{ViT-Tiny}) paradigm, we execute critical empirical proof-of-concept (PoC) audits designed to address architectural, subspace, and scale-up variations:
\begin{enumerate}
    \item \textbf{CNN Architecture (ResNet-18):} We fine-tune four expert models using a standard convolutional backbone (\texttt{resnet18}, 11.7M parameters) under identical tasks and data-scarcity constraints, and optimize merging coefficients under the source schema (\texttt{sym\_channel}) at $b=4$.
    \item \textbf{Subspace-Constrained (LoRA-like) Merging:} To mathematically replicate Parameter-Efficient Fine-Tuning (PEFT) where weight adaptations are restricted to a low-intrinsic-dimension subspace, we project our full-parameter \texttt{ViT-Tiny} expert task vectors into a low-rank subspace (rank $r=4$) via Singular Value Decomposition (SVD). We then run our standard test-time adaptation protocol on these low-rank weights.
    \item \textbf{Scale-Up Methodological Bottlenecks:} We formally examine the engineering and computational challenges that prevent direct empirical scale-up validation on multi-billion parameter backbones (e.g., \texttt{ViT-Base}, \texttt{ViT-Large}, or modern LLMs/VLMs) under this test-time paradigm, establishing a clear scientific rationale for our rigorous, controlled evaluation design.
\end{enumerate}

\begin{table*}[t]
\caption{PoC empirical results comparing the Cross-Schema Generalization Gap (Average Accuracy \%) across full-parameter ViT-Tiny, convolutional ResNet-18, and low-rank subspace-constrained (LoRA-like) ViT-Tiny. Coefficients are optimized under Symmetric Per-Channel ($Q_{\text{opt}}$) and evaluated under Symmetric Per-Tensor ($Q_{\text{eval}}$) at 4-bit precision. Note that both ResNet-18 and the Low-Rank Subspace models are trained and configured to achieve performance well above the 10.00\% random-guess floor, eliminating any potential floor-effect confounding.}
\label{tab:poc_extensions}
\vskip 0.15in
\begin{center}
\begin{small}
\begin{sc}
\begin{tabular}{lccc}
\toprule
Model / Subspace & Matched (Sym. Channel \%) & Mismatched (Sym. Tensor \%) & \textbf{Gen. Gap ($\Delta$)} \\
\midrule
ViT-Tiny (Full Parameter) & 17.88 & 10.12 & -7.76\% \\
ResNet-18 (CNN Baseline) & 21.25 & 17.00 & -4.25\% \\
ViT-Tiny (Low-Rank Subspace) & 13.00 & 13.50 & +0.50\% \\
\bottomrule
\end{tabular}
\end{sc}
\end{small}
\end{center}
\vskip -0.1in
\end{table*}

The PoC metrics summarized in Table~\ref{tab:poc_extensions} yield two highly significant discoveries:
\begin{itemize}
    \item \textbf{Architectural Mitigation:} The convolutional ResNet-18 baseline exhibits a smaller cross-schema generalization gap ($-4.25\%$) than the ViT-Tiny baseline ($-7.76\%$), while maintaining an absolute performance level (17.00\% mismatched, 21.25\% matched) significantly above the 10.00\% random-guess floor. This confirms our architectural hypothesis: localized, translation-invariant spatial kernels in CNNs have smoother weight distributions and are more resilient to the activation-shifting discretization errors caused by coarse, tensor-wise hardware target scales.
    \item \textbf{PEFT/Subspace Robustness and the ``Low-Capacity'' Generalization Illusion:} Restricting weight modifications to a low-rank subspace (which is the mathematical essence of attention-only LoRA) completely eliminates the Cross-Schema Generalization Gap ($+0.50\%$) under evaluation, with a matched accuracy of $13.00\%$. However, from a critical \emph{Methodologist} perspective, we must issue a severe caveat regarding this apparent robustness. Although the gap is closed, the absolute performance of the Low-Rank Subspace model is exceptionally low ($13.00\%$), representing a catastrophic loss of information compared to the unquantized FP16 Task Arithmetic baseline ($35.12\%$) and individual experts ($>90\%$). 

    We hypothesize that this flat generalization gap is a \textbf{Low-Capacity Generalization Illusion} rather than an active, robust alignment of expert weights. Performing a post-hoc global SVD task vector projection to compress parameter updates into a rank-4 subspace is highly destructive, stripping away vital task-specific representation directions and flattening the network's output probability distributions. Because the resulting model has highly degraded representational capacity and a small dynamic activation range, its activations exhibit minimal variance. Such a degenerate model is naturally less sensitive to quantization-operator mismatches because its predictions are already highly diffused and approach random noise under any operator. 

    Furthermore, this global, post-hoc SVD projection is a poor proxy for actual natively-trained Parameter-Efficient Fine-Tuning (such as LoRA, where $W = W_0 + BA$). Natively-trained LoRA models preserve high performance under low rank because their low-rank adapters are optimized end-to-end to align task-specific representations within a low-intrinsic dimension. Therefore, our claims of low-rank subspace restriction acting as a ``structural defense'' must be heavily tempered: the robustness observed here is partially a confounding artifact of severe capacity degradation. Evaluating actual natively-tuned LoRA experts remains a critical and necessary future extension to confirm whether the Cross-Schema Generalization Gap remains closed when high model capacity is fully preserved.
    \item \textbf{Scale-Up Methodological Bottlenecks and Analytical Projections:} We must explicitly address the physical and engineering bottlenecks that prevent running direct scale-up experiments on larger backbones (e.g., \texttt{ViT-Base} with 86M parameters or Pythia/GPT-2 language models) in this setting. In model merging research, parameter-space fusion is only mathematically coherent if task-specific experts are fine-tuned from the \emph{exact same} pre-trained weight initialization. In this workspace, pre-trained and fine-tuned expert checkpoints for the four evaluation datasets (MNIST, FashionMNIST, CIFAR-10, SVHN) are strictly provided for the \texttt{ViT-Tiny} architecture. Training a fresh set of four full-parameter experts on a larger architecture like \texttt{ViT-Base} (86M parameters) or an LLM from scratch on a multi-task sequential stream would require extensive training resources and multiple hours of execution, which is computationally and time-prohibitive under test-time constraints. 

    However, our choice of \texttt{ViT-Tiny} (5.7M parameters) represents a deliberate methodological design rather than a limitation. By maintaining complete experimental control on a lightweight architecture, we isolate the exact mathematical behaviors of the straight-through estimator and prediction entropy minimization without the compounding confounders of model capacity. As detailed in our analytical scaling projections (Section~\ref{submission:sec:conclusion}), scaling up parameter count is mathematically expected to expand, not shrink, the Cross-Schema Generalization Gap. Under unconstrained test-time adaptation, larger models introduce an exponential explosion of independent discrete rounding thresholds and dense scaling factors. This dramatically increases the risk of overfitting continuous parameters to the simulated operator's scale and zero-point parameters. Thus, the deconstruction findings on our controlled \texttt{ViT-Tiny} benchmark serve as a highly conservative lower-bound of the quantization-operator overfitting that would plague large-scale, heterogeneous hardware deployments.
\end{itemize}
